\section*{Chapitre \ref{ch:g9:thales} --- Le théorème de Thalès}

\begin{solution}{exo:g9:thales:1}
Thalès~: $\dfrac{AM}{AB} = \dfrac{AN}{AC} = \dfrac{MN}{BC}$, c'est-à-dire
$\dfrac26 = \dfrac13$. Donc $\dfrac{3}{AC} = \dfrac13$ donne $AC = 9$,
et $\dfrac{MN}{9} = \dfrac13$ donne $MN = 3$.
\end{solution}

\begin{solution}{exo:g9:thales:2}
D'abord trouver $AB = AM + MB = 5 + 3 = 8$. Puis
$\dfrac{AM}{AB} = \dfrac{AN}{AC}$ donne $\dfrac58 = \dfrac{4}{AC}$, donc
$5 \times AC = 32$ et $AC = 6.4$.
\end{solution}

\begin{solution}{exo:g9:thales:3}
Le papillon fonctionne exactement comme la configuration emboîtée~:
$\dfrac{AM}{AB} = \dfrac{AN}{AC} = \dfrac{MN}{BC}$, c'est-à-dire
$\dfrac{3}{7.5} = \dfrac25$. Donc $\dfrac{2}{AC} = \dfrac25$ donne
$AC = 5$, et $\dfrac{2.6}{BC} = \dfrac25$ donne
$BC = \dfrac{2.6 \times 5}{2} = 6.5$.
\end{solution}

\begin{solution}{exo:g9:thales:4}
$\dfrac{AM}{AB} = \dfrac68 = \dfrac34$ et
$\dfrac{AN}{AC} = \dfrac{9}{12} = \dfrac34$~: rapports égaux, points
dans le même ordre, donc $(MN) \parallel (BC)$ par la réciproque de
Thalès.
\end{solution}

\begin{solution}{exo:g9:thales:5}
$\dfrac{AM}{AB} = \dfrac46 = \dfrac23$ et
$\dfrac{AN}{AC} = \dfrac58$. Contrôle croisé~: $\dfrac23 = \dfrac{16}{24}$
et $\dfrac58 = \dfrac{15}{24}$~: les rapports diffèrent. Si les droites
étaient \omterm{def:g6:lines:perp}{parallèles}, le théorème de Thalès les forcerait à être égaux ---
donc $(MN)$ et $(BC)$ ne sont \emph{pas} \omterm{def:g6:lines:perp}{parallèles}.
\end{solution}

\begin{solution}{exo:g9:thales:6}
Soit $S$ le bout de l'ombre, $P$ le \omterm{def:g5:solids:def}{sommet} de la tête de la personne, $L$
le \omterm{def:g5:solids:def}{sommet} du lampadaire. La personne ($1.8$ m à distance $3$ m de $S$)
et le lampadaire (hauteur $h$ à distance $3 + 2 = 5$ m de $S$) sont deux
\omterm{def:g6:lines:objects}{segments} verticaux \omterm{def:g6:lines:perp}{parallèles} coupés par le \omterm{def:g3:shapes:shapes}{rayon} lumineux $(SL)$~:
Thalès depuis le point $S$ donne
\[
\frac{1.8}{h} = \frac{3}{5},
\qquad\text{donc}\quad
h = \frac{1.8 \times 5}{3} = 3 \text{ m}.
\]
\end{solution}

\begin{solution}{exo:g9:thales:7}
\emph{1.} Distance réelle~: $6.8 \times 25\,000 = 170\,000$ cm $=
1.7$ km.

\emph{2.} Les \omterm{def:g6:measure:area}{aires} s'échelonnent par $k^2$~: \omterm{def:g6:measure:area}{aire} réelle $= 8 \times
25\,000^2$ cm$^2$ $= 8 \times 6.25 \times 10^8$ cm$^2$ $= 5 \times
10^9$ cm$^2$. Puisque $1$ km$^2 = 10^{10}$ cm$^2$, cela fait $0.5$
km$^2$.
\end{solution}

\begin{solution}{exo:g9:thales:8}
\emph{1.} Le rapport est $\dfrac{AM}{AB} = \dfrac{8}{12} = \dfrac23$.
Thalès~: $AN = \frac23 \times 15 = 10$ et $MN = \frac23 \times 18 = 12$.

\emph{2.} Le \omterm{def:g9:thales:scaling}{facteur d'échelle} est $k = \frac23$. \omterm{def:g6:measure:perimeter}{Périmètres}~: $ABC$ a
$12 + 15 + 18 = 45$, et $AMN$ a $8 + 10 + 12 = 30 = \frac23 \times
45$~: le \omterm{def:g6:measure:perimeter}{périmètre} s'échelonne aussi par $k$, comme toute longueur.
\end{solution}

\begin{solution}{exo:g9:thales:9}
Autour de $O$, les droites $(AC)$ et $(BD)$ se croisent, et
$(AB) \parallel (CD)$~: configuration papillon avec les triangles $OAB$
et $OCD$. Thalès donne
\[
\frac{OA}{OC} = \frac{OB}{OD} = \frac{AB}{CD} = \frac46 = \frac23 .
\]
Donc $\dfrac{OA}{OC} = \dfrac23$, et l'égalité
$\dfrac{OA}{OC} = \dfrac{OB}{OD}$ dit précisément que $O$ coupe les deux
diagonales dans le même rapport.
\end{solution}

\begin{solution}{pb:g9:thales:1}
\textbf{1.} La construction \omterm{def:g2:mult:def}{produit} deux points sur $[AB]$~;
appelons-les $M_1$ (venant de $P_1$) et $M_2$ (venant de $P_2$).

\textbf{2.} Dans le triangle $A P_3 B$, les \omterm{def:g6:lines:perp}{parallèles} à
$(P_3 B)$ passant par $P_1$ et $P_2$ coupent les côtés
proportionnellement (\cref{thm:g9:thales:direct})~:
\[
\frac{A M_1}{AB} = \frac{A P_1}{A P_3} = \frac13,
\qquad
\frac{A M_2}{AB} = \frac{A P_2}{A P_3} = \frac23 .
\]
Le compas a rendu $A P_1$, $P_1 P_2$, $P_2 P_3$ égaux~: les rapports
sont donc exactement des tiers --- quels que soient l'inclinaison de
la demi-droite et l'écartement de compas choisis.

\textbf{3.} Pour des cinquièmes~: reporter cinq longueurs égales,
joindre le cinquième point à $B$, et tracer quatre \omterm{def:g6:lines:perp}{parallèles}. Pour
les $\frac35$ d'un \omterm{def:g6:lines:objects}{segment}~: même figure, et prendre le point
découpé par la \omterm{def:g6:lines:perp}{parallèle} issue de $P_3$~: il se trouve aux
$\frac{3}{5}$ de $[AB]$ depuis $A$.

\textbf{4.} $\frac{AM}{MB} = \frac23$ signifie
$AM = \frac25 AB$~: cinq reports égaux sur la demi-droite, joindre
$P_5$ à $B$, et la \omterm{def:g6:lines:perp}{parallèle} issue de $P_2$ marque $M$.

\textbf{5.} Mesurer $\sqrt2 = 1.41421\ldots$ cm ne mobilise jamais
qu'un nombre fini de décimales~: tout «~tiers~» mesuré n'est qu'une
approximation. La construction de Thalès, elle, ne lit aucun
nombre~: elle livre le point exact situé à $\frac{\sqrt2}{3}$ --- des
parts égales d'un \omterm{def:g6:lines:objects}{segment} qu'aucune règle ne sait même exprimer
(\cref{pb:g9:sqrt:1}).

\textbf{6.} Le \omterm{def:g5:solids:def}{sommet} de l'objet, le trou et le bas de l'image sont
alignés, et de même pour le bas de l'objet et le \omterm{def:g5:solids:def}{sommet} de l'image~:
deux droites se coupant au trou, avec l'objet et la paroi de l'image
\omterm{def:g6:lines:perp}{parallèles}. Thalès dans le papillon~:
\[
\frac{h}{H} = \frac{d}{D},
\qquad\text{donc}\qquad
h = H \times \frac dD .
\]

\textbf{7.} $h = 12 \times \frac{0.20}{30} = 0.08$ m $= 8$ cm --- à
l'envers (les \omterm{def:g3:shapes:shapes}{rayons} se croisent au trou).

\textbf{8.} $H = h \times \frac Dd = 9 \times 10^{-3} \times
\frac{1.5 \times 10^{11}}{1} = 1.35 \times 10^{9}$ m. Valeur
réelle~: $1.39 \times 10^9$ m~: une boîte à chaussures mesure le
Soleil à quelques pour cent près.

\textbf{9.} Lune~: $\frac{3.5 \times 10^6}{3.8 \times 10^8}
\approx 0.0092$. Soleil~: $\frac{1.39 \times 10^9}{1.5 \times
10^{11}} \approx 0.0093$. Les deux rapports --- les tailles
apparentes dans le ciel --- sont presque identiques~: la Lune peut
recouvrir le Soleil \emph{exactement}, bord à bord. Ce hasard
cosmique, c'est l'éclipse totale de Soleil.

\textbf{10.} Tout \omterm{def:g3:shapes:shapes}{rayon} doit passer par l'unique trou~: les \omterm{def:g3:shapes:shapes}{rayons}
venant du haut de l'objet poursuivent donc \emph{vers le bas} et
ceux venant du bas poursuivent \emph{vers le haut}~: l'image est
renversée. Agrandir le trou laisse passer tout un faisceau de copies
légèrement décalées de l'image, qui se superposent~: plus lumineuse,
mais brouillée.

\textbf{11.} $\frac{PQ}{MN} = \frac{OP}{OM} = \frac{7.5}{3}
= 2.5$, donc $PQ = 10$~; et $\frac{ON}{OQ} = \frac{OM}{OP} =
\frac{1}{2.5}$, donc $ON = \frac{6}{2.5} = 2.4$.

\textbf{12.} Si la \omterm{def:g6:lines:perp}{parallèle} passe par le milieu d'un côté, le
rapport de Thalès vaut $\frac12$~: elle coupe donc le second côté en
\emph{son} milieu, et le \omterm{def:g6:lines:objects}{segment} \omterm{def:g6:lines:perp}{parallèle} mesure la \omterm{def:g3:division:half}{moitié} de la
base --- c'est précisément \cref{thm:g8:midpoints:direct}. Thalès
est le théorème des milieux affranchi du rapport $\frac12$.

\textbf{13.} Dans le triangle $ACD$, la droite $(EO)$ est \omterm{def:g6:lines:perp}{parallèle}
à la base $(CD)$, avec $\frac{AO}{AC} = \frac{OA}{OA + OC} =
\frac{4}{4 + 6} = \frac25$. Thalès~:
$\frac{EO}{DC} = \frac{AO}{AC} = \frac25$, donc
$EO = \frac25 \times 6 = 2.4$.

\textbf{14.} Dans le triangle $BDC$, de même,
$\frac{BO}{BD} = \frac25$ et $OF = \frac25 \times 6 = 2.4$. D'où
\[
EF = EO + OF = 4.8
= \frac{2 \times 4 \times 6}{4 + 6} :
\]
la \omterm{pb:g8:speed:1}{moyenne harmonique} des bases --- la moyenne des aller-retours
(\cref{pb:g8:speed:1}), dessinée dans un trapèze.

\textbf{15.} Arithmétique~: $5$~; géométrique~:
$\sqrt{24} \approx 4.90$~; harmonique~: $4.8$. Rangement~:
$4.8 < 4.90\ldots < 5$, c'est-à-dire harmonique $<$ géométrique $<$
arithmétique --- avec égalité seulement pour des bases égales. Dans
la figure~: la ligne des milieux mesure $5$, la \omterm{def:g6:lines:perp}{parallèle} passant
par le croisement des diagonales $4.8$, et la coupe de similitude
$\sqrt{24}$, toutes empilées entre les bases dans cet ordre. Les
inégalités ont été démontrées dans \cref{pb:g8:cosine:1}
(géométrique $\leq$ arithmétique, par le cercle et par une identité)
et dans \cref{pb:g8:speed:1} (harmonique $\leq$ arithmétique)~;
harmonique $\leq$ géométrique s'en déduit en les combinant
($H \times A = G^2$, démontré là aussi).
\end{solution}
